% GENERATED FILE. Edit canonical lessons, then run npm run generate:books.
% canonical-source-hash: fnv1a64:b8dcaeb542ee5a3e
% canonical-lessons: ZH-C334-zhengyue, ZH-C334-hongbao, ZH-W334-hong-red, ZH-C334-tie, ZH-C334-chunlian, ZH-C334-honghuo

\chapter{Red Envelopes}
\label{ch:zh-red-envelopes}
\hlchaptermodality{334}

\begin{chapteropening}
\textbf{By the end of this chapter:} \emph{I can talk about the first lunar month, red envelopes, couplets by the door and a business doing well.}

Complete the last lesson of chapter 334: I can talk about the first lunar month, red envelopes, couplets by the door and a business doing well.
\end{chapteropening}

\section[\texorpdfstring{\zh{正月} (zhēngyuè)}{正月 (zhēngyuè)}]{\zh{正月} (zhēngyuè) — first lunar month}

\label{lesson:ZH-C334-zhengyue}



\begin{quote}
Before the new one: say the Mandarin for to keep the New Year vigil, then the Mandarin for to pay a New Year call.
\end{quote}

\subsection*{You'll want to know: \zh{正月}}
\textbf{\zh{正月}} — \emph{zhēngyuè} — \textquotedblleft{}first lunar month\textquotedblright{}.

\begin{quote}
\textbf{\zh{正月里大家都很开心}} \emph{zhēngyuè lǐ dàjiā dōu hěn kāixīn} — In the first month everyone is cheerful.
\end{quote}

\zh{正} is said zhēng here, with a level tone, not the zhèng of \zh{正在}.

\begin{cousinweb}
\textbf{\zh{正月}} (\emph{zhēngyuè}) is built from \zh{正} and \zh{月} 'a month', both inherited from Middle Chinese, \emph{tsyeng} and \emph{ngjwot} (Guangyun, in Baxter's transcription). The dictionary gives the first its level-tone reading \emph{tsyeng} for the start of the calendar year, apart from \emph{tsyengH}, 'upright': so zhēng here, and zhèng in \zh{正在}.
\end{cousinweb}

\begin{grammarlens}[title={What you've built}]
One more word for telling what happened.
\end{grammarlens}

\subsection*{Your turn}
\begin{itemize}
\raggedright
  \item \emph{Say it:} \emph{zhēngyuè}
  \item \emph{Say it:} \emph{zhēngyuè}, once more
  \item \emph{Say it:} \emph{zhēngyuè lǐ dàjiā dōu hěn kāixīn} — In the first month everyone is cheerful.
  \item \emph{Recall:} say the Mandarin for to keep the New Year vigil, then the Mandarin for to pay a New Year call, then say \emph{zhēngyuè} again
\end{itemize}

\begin{tcolorbox}[breakable,colback=teal!4,colframe=teal!35!black,title={Before you move on}]
What does \zh{正月} mean? (\textquotedblleft{}First lunar month\textquotedblright{}.) Say it once more.
\end{tcolorbox}
\section[\texorpdfstring{\zh{红包} (hóngbāo)}{红包 (hóngbāo)}]{\zh{红包} (hóngbāo) — red envelope}

\label{lesson:ZH-C334-hongbao}



\begin{quote}
Before the new one: say the Mandarin for to pay a New Year call, then the Mandarin for the first lunar month.
\end{quote}

\subsection*{You'll want to know: \zh{红包}}
\textbf{\zh{红包}} — \emph{hóngbāo} — \textquotedblleft{}red envelope\textquotedblright{}.

\begin{quote}
\textbf{\zh{过年的时候外公给了我一个红包}} \emph{guònián de shíhou wàigōng gěi le wǒ yí ge hóngbāo} — At the new year my grandpa gave me a red envelope.
\end{quote}

The new character, in pieces: \zh{红} — six strokes: on the left the silk piece of \zh{给}, three strokes; on the right \zh{工}. The next lesson puts it in your hand.

\begin{cousinweb}
\textbf{\zh{红包}} (\emph{hóngbāo}) is built from \zh{红} 'red' and \zh{包} 'a packet', both inherited from Middle Chinese, \emph{huwng} and \emph{paew} (Guangyun, in Baxter's transcription), which glosses the first 'a colour' and the second 'to wrap'.
\end{cousinweb}

\begin{grammarlens}[title={What you've built}]
One more word for telling what happened.
\end{grammarlens}

\subsection*{Your turn}
\begin{itemize}
\raggedright
  \item \emph{Say it:} \emph{hóngbāo}
  \item \emph{Say it:} \emph{hóngbāo}, once more
  \item \emph{Say it:} \emph{guònián de shíhou wàigōng gěi le wǒ yí ge hóngbāo} — At the new year my grandpa gave me a red envelope.
  \item \emph{Recall:} say the Mandarin for to pay a New Year call, then the Mandarin for the first lunar month, then say \emph{hóngbāo} again
\end{itemize}

\begin{tcolorbox}[breakable,colback=teal!4,colframe=teal!35!black,title={Before you move on}]
What does \zh{红包} mean? (\textquotedblleft{}Red envelope\textquotedblright{}.) Say it once more.
\end{tcolorbox}
\section[\texorpdfstring{\zh{红} (hóng)}{红 (hóng)}]{\zh{红} — the silk piece, then \zh{工}}

\label{lesson:ZH-W334-hong-red}



\begin{quote}
The word first, then the character inside it.

\begin{itemize}
\raggedright
  \item \emph{Recall:} say \emph{hóngbāo} — \textbf{R1}, one lesson back, before it decays
\end{itemize}

\begin{quote}
\zh{给} \zh{工}
\end{quote}

\emph{Read it:} each one aloud, with nothing said first — \emph{gěi}, \emph{gōng}
\end{quote}

\subsection*{Script — the silk piece, then \zh{工}}
\hlblockfigure{\detokenize{figures/ZH-W334-hong-red-filmstrip.pdf}}{How \zh{红} is written, stroke by stroke}

\begin{quote}
\zh{红}
\end{quote}

\textbf{Two pieces, left before right.}

\textbf{The left piece, three strokes,} as on the left of \zh{给}: a stroke that falls to the left and turns right, a second one like it below, then a short stroke rising to the right.

\textbf{The right piece, \zh{工}, three strokes.} A short bar, a vertical down from its middle, then a longer bar along the foot.

3 + 3 = six.

\subsection*{Your turn}
\begin{itemize}
\raggedright
  \item \emph{Write it:} \zh{红}, all six, with the model in front of you
  \item \emph{Cover:} write it again from memory, and count the pieces, not the strokes
\end{itemize}

\begin{tcolorbox}[breakable,colback=teal!4,colframe=teal!35!black,title={Before you move on}]
How many strokes in \zh{红}? (\textbf{Six}.) Where have you met the left piece? (\textbf{On the left of \zh{给}}.) What is the right piece? (\textbf{\zh{工}}.)

Source: \href{https://raw.githubusercontent.com/chanind/hanzi-writer-data/68d10a4b21150cae5e1ebbd223eed289cf32d90c/data/\%E7\%BA\%A2.json}{Hanzi Writer Data \zh{红}, pinned stroke paths and medians}.
\end{tcolorbox}
\section[tiē]{tiē — to stick; to paste up}

\label{lesson:ZH-C334-tie}



\begin{quote}
Before the new one: say the Mandarin for the first lunar month, then the Mandarin for a red envelope of gift money.
\end{quote}

\subsection*{You'll want to know: tiē}
\emph{tiē} — \textquotedblleft{}to stick; to paste up\textquotedblright{}.

\begin{quote}
\emph{wǒ zài ménkǒu tiē le yì zhāng hóng zhǐ} — I stuck a sheet of red paper up by the door.
\end{quote}

Its characters come later; for now it is a word for the ear.

\begin{cousinweb}
\emph{tiē} is inherited from Middle Chinese \emph{thep} (Guangyun, in Baxter's transcription), which glosses it 'to pawn something for money'. The sense 'to stick on' is not given there, and how it came is not settled.
\end{cousinweb}

\begin{grammarlens}[title={What you've built}]
One more word for telling what happened.
\end{grammarlens}

\subsection*{Your turn}
\begin{itemize}
\raggedright
  \item \emph{Say it:} \emph{tiē}
  \item \emph{Say it:} \emph{tiē}, once more
  \item \emph{Say it:} \emph{wǒ zài ménkǒu tiē le yì zhāng hóng zhǐ} — I stuck a sheet of red paper up by the door.
  \item \emph{Recall:} say the Mandarin for the first lunar month, then the Mandarin for a red envelope of gift money, then say \emph{tiē} again
\end{itemize}

\begin{tcolorbox}[breakable,colback=teal!4,colframe=teal!35!black,title={Before you move on}]
What does tiē mean? (\textquotedblleft{}To stick; to paste up\textquotedblright{}.) Say it once more.
\end{tcolorbox}
\section[chūnlián]{chūnlián — Spring Festival couplets}

\label{lesson:ZH-C334-chunlian}



\begin{quote}
Before the new one: say the Mandarin for a red envelope of gift money, then the Mandarin for to paste up on a door.
\end{quote}

\subsection*{You'll want to know: chūnlián}
\emph{chūnlián} — \textquotedblleft{}Spring Festival couplets\textquotedblright{}.

\begin{quote}
\emph{Chúxī de xiàwǔ wǒmen zài ménkǒu tiē le chūnlián} — On the afternoon of the last day of the year we put up the couplets by the door.
\end{quote}

Its characters come later; for now it is a word for the ear.

\begin{cousinweb}
\emph{chūnlián} is built from \zh{春} 'spring' and \emph{lián} 'to join'. The first is inherited from Middle Chinese \emph{tsyhwin} (Guangyun, in Baxter's transcription), the first of the four seasons. The dictionary lists the second only under an older spelling of the character, \emph{ljen}, which it glosses 'unbroken'.
\end{cousinweb}

\begin{grammarlens}[title={What you've built}]
One more word for telling what happened.
\end{grammarlens}

\subsection*{Your turn}
\begin{itemize}
\raggedright
  \item \emph{Say it:} \emph{chūnlián}
  \item \emph{Say it:} \emph{chūnlián}, once more
  \item \emph{Say it:} \emph{Chúxī de xiàwǔ wǒmen zài ménkǒu tiē le chūnlián} — On the afternoon of the last day of the year we put up the couplets by the door.
  \item \emph{Recall:} say the Mandarin for a red envelope of gift money, then the Mandarin for to paste up on a door, then say \emph{chūnlián} again
\end{itemize}

\begin{tcolorbox}[breakable,colback=teal!4,colframe=teal!35!black,title={Before you move on}]
What does chūnlián mean? (\textquotedblleft{}Spring Festival couplets\textquotedblright{}.) Say it once more.
\end{tcolorbox}
\section[\texorpdfstring{\zh{红火} (hónghuǒ)}{红火 (hónghuǒ)}]{\zh{红火} (hónghuǒ) — prosperous; flourishing}

\label{lesson:ZH-C334-honghuo}



\begin{quote}
Before the new one: say the Mandarin for to paste up on a door, then the Mandarin for the Spring Festival couplets.
\end{quote}

\subsection*{You'll want to know: \zh{红火}}
\textbf{\zh{红火}} — \emph{hónghuǒ} — \textquotedblleft{}prosperous; flourishing\textquotedblright{}.

\begin{quote}
\textbf{\zh{这家饭店很红火}} \emph{zhè jiā fàndiàn hěn hónghuǒ} — This restaurant is doing very well.
\end{quote}

\begin{cousinweb}
\textbf{\zh{红火}} (\emph{hónghuǒ}) is built from \zh{红} 'red' and \zh{火} 'fire', both inherited from Middle Chinese, \emph{huwng} and \emph{xwaX} (Guangyun, in Baxter's transcription), which glosses the first 'a colour' and the second, quoting an older dictionary, 'to blaze'.
\end{cousinweb}

\begin{grammarlens}[title={What you've built}]
One more word for telling what happened.
\end{grammarlens}

\subsection*{Your turn}
\begin{itemize}
\raggedright
  \item \emph{Say it:} \emph{hónghuǒ}
  \item \emph{Say it:} \emph{hónghuǒ}, once more
  \item \emph{Say it:} \emph{zhè jiā fàndiàn hěn hónghuǒ} — This restaurant is doing very well.
  \item \emph{Recall:} say the Mandarin for to paste up on a door, then the Mandarin for the Spring Festival couplets, then say \emph{hónghuǒ} again
\end{itemize}

\begin{tcolorbox}[breakable,colback=teal!4,colframe=teal!35!black,title={Before you move on}]
What does \zh{红火} mean? (\textquotedblleft{}Prosperous; flourishing\textquotedblright{}.) Say it once more.
\end{tcolorbox}
